\section*{Chapter \ref{ch:b3:environmental-and-resource-economics} --- Environmental and Resource Economics}

\begin{solution}{exo:b3:environmental-and-resource-economics:1}
The United Kingdom in 1907, 134 kg per person, at 7{,}855 dollars; China in 2006, 29 kg, at 7{,}271 dollars --- similar incomes, a
century apart, and peaks a fifth as high. % figdata: so2_stats
\end{solution}

\begin{solution}{exo:b3:environmental-and-resource-economics:2}
$2x_1 = 3x_2$ and $x_1 + x_2 = 100$: $x_1 = 60$, $x_2 = 40$, and a tax of 120 dollars a tonne.
\end{solution}

\begin{solution}{exo:b3:environmental-and-resource-economics:3}
$50 \times 1.03^{10} = 67.20$ dollars a tonne.
\end{solution}

\begin{solution}{exo:b3:environmental-and-resource-economics:4}
Real net prices flat or falling over decades. Before concluding, check that
extraction costs and known stocks did not change: falling costs or new discoveries
lower the path without contradicting the arbitrage.
\end{solution}

\begin{solution}{exo:b3:environmental-and-resource-economics:5}
$1/r$: 71.4, 33.3 and 20.0 dollars.
\end{solution}

\begin{solution}{exo:b3:environmental-and-resource-economics:6}
$0.1 + 1 \times 1.3 = 1.4~\%$ and $1.5 + 2 \times 1.3 = 4.1~\%$.
\end{solution}

\begin{solution}{exo:b3:environmental-and-resource-economics:7}
The tested solution is \code{ex07_so2_peaks.py}, in the chapter's \code{exercises/}
folder. France peaked in 1976 at 58.6 kg per person and 21{,}465 dollars; Italy in 1980 at 59.1 kg and 20{,}959 dollars; Spain in 1981 at 77.3 kg
and 13{,}895 dollars. In 2019 they emitted 1.5, 1.7 and 3.7 kg. % output: ex07_so2_peaks
\end{solution}

\begin{solution}{exo:b3:environmental-and-resource-economics:8}
$100 \times (8 - 2)/(2 \times 64) = 4.69$: the tax is better, as for carbon dioxide, whose damage depends on the accumulated
stock and hardly on one year's emissions. $100 \times (8 - 20)/128 = -9.38$: the cap is better, as for a pollutant with a
threshold of harm.
\end{solution}

\begin{solution}{exo:b3:environmental-and-resource-economics:9}
Falls that start at the dates of laws, at different incomes in different
countries; countries at the same income with different rules emitting very
different amounts; and within a country, the fall concentrated in regulated sources.
\end{solution}

\begin{solution}{exo:b3:environmental-and-resource-economics:10}
Extracting at $t$ and investing gives $(P_t - c)(1 + r\,dt)$ at $t + dt$; waiting gives $P_{t+dt} - c$. Indifference gives $P_{t+dt} - P_t = r(P_t - c)\,dt$, so
$\dot P/(P - c) = r$ and $P_t = c + (P_0 - c)e^{rt}$.
\end{solution}

\begin{solution}{exo:b3:environmental-and-resource-economics:11}
$10 = 1{,}000^{b}$, so $b = 1/3$; the \omterm{def:b3:environmental-and-resource-economics:learning-rate}{learning rate} is $1 - 2^{-1/3} = 20.6~\%$.
\end{solution}

\begin{solution}{exo:b3:environmental-and-resource-economics:12}
The tested solution is \code{ex12_scc.py}, in the chapter's \code{exercises/}
folder. 1{,}217.6 dollars at 1.4~\%, 142.0 at 3~\% and 50.0 at 5~\%; 508.4 with the declining rate. At 1.4~\% the damages grow faster than
the discount rate and the sum depends on where it is cut off. % output: ex12_scc
\end{solution}

\begin{solution}{pb:b3:environmental-and-resource-economics:coal}
% figdata: coal_stats
\textbf{1.} $P_t = 20 + \lambda e^{0.04t}$; $\lambda$ is such that $\int_0^T (40 - 0.25P_t)\,dt = 600$, where $P_T = 100$.

\textbf{2.} $\lambda = 27.75$ dalens; $P = 47.75$.

\textbf{3.} 28.06 million tonnes.

\textbf{4.} Rent 41.39, price 61.39, extraction 24.65 million tonnes.

\textbf{5.} After 26.5 years, in 2051.

\textbf{6.} A lower rent today and a later date of exhaustion.

\textbf{7.} Prices are capped lower, so more is extracted each year and the coal runs
out sooner, at a lower rent.

\textbf{8.} $28.06 \times 0.02 = 0.561$ million tonnes, 561 kilotonnes.

\textbf{9.} $561 \times 0.4 = 224.5$.

\textbf{10.} $0.4 \times 3.0 \text{ million}/1{,}000 = 1{,}200$ dalens a tonne.

\textbf{11.} 673.5 million dalens a year.

\textbf{12.} The dose--response relation, from epidemiology; the \omterm{def:b3:cost-benefit-analysis:vsl}{value of a statistical life}
comes second.

\textbf{13.} The smoke is a cost of burning coal; once priced, demand for coal falls at
every price and the coal lasts longer.

\textbf{14.} The marginal damage, 1{,}200 dalens a tonne.

\textbf{15.} $x_1 = 1{,}200/10 = 120$ and $x_2 = 1{,}200/30 = 40$ kilotonnes.

\textbf{16.} $\tfrac12 \times 10 \times 120^2 + \tfrac12 \times 30 \times 40^2 = 96$ million dalens.

\textbf{17.} 80 each: $\tfrac12 \times 40 \times 80^2 = 128$ million; the tax saves 32 million.

\textbf{18.} A cap of $561.3 - 160 = 401.3$ kilotonnes; permits trade at 1{,}200 dalens.

\textbf{19.} $160 \times 0.4 = 64$ deaths, worth $64 \times 3.0 = 192$ million dalens, twice the cost.

\textbf{20.} Standards are easier to monitor, avoid local hot spots near a plant that
buys permits, and are familiar to firms and courts.

\textbf{21.} $40 \times 190 = 7{,}600$ million dalens.

\textbf{22.} $40 \times 120 = 4{,}800$ and $40 \times 340 = 13{,}600$ million.

\textbf{23.} $7{,}600/393{,}264 = 1.9~\%$.

\textbf{24.} For example $\rho = 0.7~\%$ with $\eta = 1$, or $\rho = 0.05~\%$ with $\eta = 1.5$.

\textbf{25.} Most of the damage from Brenland's emissions falls on other countries; a tax
set on its own damage alone would be far lower than the social cost.

\textbf{26.} The efficient tax is 1{,}200 dalens a tonne of sulphur dioxide, and Brenland's coal is exhausted in 2051.
\end{solution}
